\subsection{Linear quasi-potential bath model}


Following methodology well established in the literature \cite{milewski2015faraday,galeano2017non,galeano2021capillary, alventosa2023inertio} to describe the waves generated by a bath-droplet interaction, we proceed by assuming a fluid bath with small free-surface perturbations from the equilibrium $z=0$  denoted by $z=\eta_b (x,t)$. As the liquid is incompressible we can take the velocity field of the bath as  $\mathbf{u}_b = (u_b, w_b) = \nabla \phi + \nabla^\perp \psi$, where $\nabla \phi$ denotes the irrotational velocity component, and $\psi$ a small vortical correction arising from viscous effects. Here we have used the notation $\nabla^\perp = (-\partial_z, \partial_x)$. Due to the assumption that disturbances are small, the system \eqref{eq: mass} -- \eqref{eq:VelocityCont} results in the following linearized problem

\begin{align}
\nabla^2 \phi &= 0, \qquad &z&\le 0, \label{eq: laplace}\\
\rho_b \partial_t \phi &= - p_b - \rho_b g z, \qquad &z&\le 0, \\
\rho_b \partial_t \psi &= \mu_b \nabla^2 \psi, \qquad &z& \le 0, 
\end{align}
and boundary conditions

\begin{align}
\partial_t \eta_b  &= \partial_z \phi + \partial_x \psi , \qquad &z&=0,  \\
 p_b - p_a - \sigma_b \kappa_b &= 2\mu_b\left(\partial_z^2 \phi + \partial_{xz} \psi \right), \qquad &z&=0, \\
   \mu_a (\partial_z u_a + \partial_x w_a) &= \mu_b (2\partial_{xz} \phi + \partial_x^2 \psi - \partial_z^2 \psi) , \qquad &z&=0, \\
   u_a &= \partial_{x} \phi - \partial_z \psi, \qquad &z&=0, \\
    w_a &= \partial_{z} \phi + \partial_x \psi, \qquad &z&=0, \label{eq: z velocity match} \\
\phi,\psi &\to 0, \qquad &z &\to -\infty, \label{eq: end of full system}
\end{align}
which include the kinematic boundary condition, normal and tangential stress balances, velocity continuity, and far field decay, respectively. We have omitted normal stresses in the air, in anticipation of a lubrication layer approximation.

The arguments in Dias \textit{et al} \cite{dias2008theory} (building on the viscous damping rates for water waves obtained by Lamb \cite{lamb1924hydrodynamics}) involve a smallness assumption in the boundary layer flow represented by $\psi$ and a boundary layer scaling for $\psi$ in its vertical dependence. This is explained in more detail in Appendix \ref{App: Boundary}. The boundary conditions then simplify to

\begin{align}
\partial_t \eta_b  &= \partial_z \phi + \partial_x \psi , \qquad &z&=0,  \\
\partial_t \phi &=  - \frac{1}{\rho_b} p_a + g \eta_b + \frac{\sigma_b}{\rho_b} \partial_x^2 \eta_{b} - 2\nu_b \partial_z^2 \phi, \qquad &z&=0, \\
    \partial_t \psi &=  \nu_b \left(2+\frac{\mu_a}{ \mu_b}\right) \partial_{zx} \phi + \nu_b \frac{\mu_a}{\mu_b}\partial_z u_a , \qquad &z&=0, \\
   u_a &= \partial_{x} \phi, \qquad &z&=0, \\
    w_a &= \partial_{z} \phi, \qquad &z&=0,
\end{align}
where we have introduced  the kinematic viscosity $\nu_b = \mu_b/\rho_b$ and substituted for the linearized curvature of the surface $\kappa = -\partial_x^2 \eta_b$. Disregarding terms of $O({\mu_a}/{\mu_b})$, one can deduce using the leading order balance $\partial_t \eta_b = \partial_z \phi$ that $\partial_x \psi = 2\nu_b \partial_{x}^2 \eta_b$, and hence the boundary conditions for the velocity potential and the free surface evolution are forced by air pressure alone

\begin{align}
\partial_t \phi &= -\frac{1}{\rho_b} p_a -  g\eta_b + \frac{\sigma_b}{\rho_b}\partial_x^2 \eta_b  + 2\nu_b \partial_x^2 \phi, \qquad &z&=0, \label{eq: qp DBC} \\
\partial_t \eta_b  &= \partial_z \phi + 2\nu_b \partial_x^2 \eta_b, \qquad &z&=0, \label{eq: qp KBC}
\end{align}
where $\phi$ satisfies Laplace's equation in the lower half plane. The conditions for velocity continuity at the interface will be matched in the equations for the lubrication air layer, and the pressure $p_a$ will be taken to be zero outside this lubrication layer. We note that the kinematic boundary condition requires $\partial_z \phi$ to be evaluated at $z=0$, and that this can be expressed simply in terms of $\phi$ at the surface once the system is Fourier transformed in $x$ as shown in Section \ref{sect:Numerics}.